% gz-p8-05-2: 医学检验贝叶斯反差柱状图——人群分布 / 阳性来源 / 阳性后患病概率
\documentclass[tikz,border=4pt]{standalone}
\usepackage{ctex}
\usepackage{amsmath}
\usetikzlibrary{calc}

\begin{document}
\begin{tikzpicture}[scale=1.0, thick]
  % 三柱柱状图，每柱宽 1.5，间距 2

  % 第 1 柱：人群分布（1% 患病 / 99% 健康）—— 高度 = 5
  \def\bx{0}
  \def\bh{5}
  % 健康（99%）
  \fill[blue!30] (\bx, 0) rectangle (\bx + 1.5, \bh * 0.99);
  % 患病（1%）
  \fill[red!50] (\bx, \bh * 0.99) rectangle (\bx + 1.5, \bh);
  % 边框
  \draw (\bx, 0) rectangle (\bx + 1.5, \bh);
  % 标签
  \node[font=\footnotesize, blue] at (\bx + 0.75, \bh * 0.5) {健康 99\%};
  \node[font=\tiny, red, above] at (\bx + 0.75, \bh) {患病 1\%};
  \node[below, font=\small] at (\bx + 0.75, -0.3) {人群分布};

  % 第 2 柱：阳性来源 - 总阳性 = 0.0095 + 0.099 = 0.1085
  % 真阳 0.0095 / 0.1085 ≈ 8.76%
  % 假阳 0.099 / 0.1085 ≈ 91.24%
  \def\bx{4}
  % 假阳（来自健康）91.24%
  \fill[blue!30] (\bx, 0) rectangle (\bx + 1.5, \bh * 0.9124);
  % 真阳（来自患病）8.76%
  \fill[red!50] (\bx, \bh * 0.9124) rectangle (\bx + 1.5, \bh);
  \draw (\bx, 0) rectangle (\bx + 1.5, \bh);
  \node[font=\footnotesize, blue, align=center] at (\bx + 0.75, \bh * 0.45)
    {假阳\\0.099};
  \node[font=\tiny, red, above] at (\bx + 0.75, \bh) {真阳 0.0095};
  \node[below, font=\small, align=center] at (\bx + 0.75, -0.3) {阳性来源};

  % 第 3 柱：阳性后患病概率（8.76% 患病 / 91.24% 健康）
  \def\bx{8}
  % 健康（91.24%）
  \fill[blue!30] (\bx, 0) rectangle (\bx + 1.5, \bh * 0.9124);
  % 患病（8.76%）
  \fill[red!50] (\bx, \bh * 0.9124) rectangle (\bx + 1.5, \bh);
  \draw (\bx, 0) rectangle (\bx + 1.5, \bh);
  \node[font=\footnotesize, blue, align=center] at (\bx + 0.75, \bh * 0.45)
    {健康\\91.24\%};
  \node[font=\tiny, red, above] at (\bx + 0.75, \bh) {患病 8.76\%};
  \node[below, font=\small, align=center] at (\bx + 0.75, -0.3) {阳性后\\$P(\text{患病}\mid+)$};

  % 总标题/箭头
  \draw[-stealth, thick, gray] (1.7, 4) -- (3.9, 4);
  \draw[-stealth, thick, gray] (5.7, 4) -- (7.9, 4);

  % 注释
  \node[font=\footnotesize, align=center] at (5, -1.4)
    {反差: 患病率从 1\% 升到 8.76\%, 但绝大多数阳性仍来自健康人};
\end{tikzpicture}
\end{document}
